\chapter{Конструкторская часть}
%<пишем про алгоритмы, архитектуру ПО и тп>

\section{Описание алгоритмов}

\begin{figure}[!htbp]
	\centering
	\includesvg[width=0.2\textwidth]{aa_lab1-std}
	\caption{Cхема стандартного алгоритма}
\end{figure}
\clearpage

\begin{figure}[!htbp]
	\centering
	\includesvg[width=0.7\textwidth]{aa_lab1-win}
	\caption{Cхема алгоритма Винограда}
\end{figure}
\clearpage

\begin{figure}[!htbp]
	\centering
	\includesvg[width=0.4\textwidth]{aa_lab1-win2}
	\caption{Cхема алгоритма Винограда}
\end{figure}
\clearpage

\begin{figure}[!htbp]
	\centering
	\includesvg[width=0.8\textwidth]{aa_lab1-win_opt}
	\caption{Cхема алгоритма Винограда с оптимизацией}
\end{figure}
\FloatBarrier


\section{Расчёт трудоёмкости алгоритмов}

Введём модель вычислений для определения теоретической трудоёмкости описанных ранее алгоритмов.\\

\begin{table}[!htbp]
	\caption{Цены базовых операций.}
	\centering
	\begin{tabularx}{\textwidth}{|*{2}{>{\centering\arraybackslash}X|}}
		\hline
		\textbf{Операции} & \textbf{Цена} \\\hline
		{ +, - , =, ++, $--$ , <<, >> } & 1 \\\hline
		{ *, /, \%, *=, /= } & 2 \\\hline
		{ ==, >=, <=, !=, >, <, \&\&, || } & 1 \\\hline
		{ [], \{\} } & 1\\\hline
	\end{tabularx}
\end{table}

\begin{table}[!htbp]
	\centering
	\caption{Трудоёмкость алгоритмов.}
	\small
	\begin{tabularx}{\textwidth}{|*{3}{X|}}
		\hline
		\textbf{Алгоритм} & \textbf{Трудоёмкость ЛС} & \textbf{Трудоёмкость ХС} \\\hline
		Стандартный & $2 + 2N + 2NQ + 10MNQ$ & $2 + 2N + 2NQ + 10MNQ$ \\\hline
		Виноград & $12 - 7N - 9Q + 6NM + 9NQ + 13.5NMQ$ & $14 - 5N - 9Q + 6NM + 18NQ + 13.5NMQ$ \\\hline
		Виноград с оптимизацией & $10 + Q + 2.25MQ + 2N + 7NQ + 13.5NMQ$ & $13 + Q + 2.25MQ + 4N + 16NQ + 13.5NMQ$ \\\hline
	\end{tabularx}
\end{table}
\FloatBarrier

\section*{Вывод}
%<что сделали в конструкторке и получили в результате, кратко>

В данном разделе мной были построены схемы алгоритмов и вычислена их трудоёмкость с помощью введённой модели вычислений.

\clearpage
